% 원제: 요약
\chapter*{요약}\addcontentsline{toc}{chapter}{요약}\markboth{요약}{요약}

PRISM-Hub는 800 V급 배터리 차량에서 대전력은 2:1 고정비 공진형 스위치드 커패시터 변환기(Stage A)가 전달하고, 조정은 소용량 부분전력 조정기(Stage B)가 맡도록 역할을 나눈 통합 직류 전력 허브다. 400·500 V급 충전기는 Stage A 출력 노드 V\textsubscript{A}에 접속하므로 충전 대전류는 조정단을 지나지 않는다. 같은 하드웨어가 주행 중에는 400 V 보조버스와 48 V 공급원이 된다.

2026년 10월 22일에 기준 구조와의 손실·정격 비교에서 Stage B의 400 V 전 범위 조정이 조정 없는 구조보다 손실·반도체 수·스위치 VA가 모두 크다는 결과가 나와(가설 H5 기각), 아키텍처를 \textbf{개선 구조 R2}로 고쳤다.

\begin{itemize}
\item Stage A를 충전 호환·400 V 중간버스·48 V 공급의 핵심으로 둔다. 기본형의 400 V 버스는 V\textsubscript{A}를 따른다.
\item Stage B는 V\textsubscript{A}가 400 V급 부품 상한을 넘는 팩에만 붙이는 1–5 kW \textbf{창 조정 모듈}로 줄이고, 바이패스로 우회할 수 있게 한다.
\item Stage C는 입력 V\textsubscript{A} 범위(270–504 V)를 받는 조정형으로 바꾼다.
\end{itemize}

창 조정 모듈의 축소 정격·모드 전환·보호 상태기계·고장 강등은 10월 25일 보고에서 평균모델과 회로 모델로 확인했다(6–9장). 이번 보고에서는 R2를 기본형까지 완성해 실제 세 단을 한 회로로 확인하고, 기존 산업 방식 대비 위치를 같은 손실 규칙으로 정량화했다. 추가한 검증은 다음과 같다.

\begin{itemize}
\item \textbf{조정형 Stage C(E22, TB14·TB15, 10.1–10.2절).} 동기 벅 + 공진점 CLLC 직류변압기(4.8:1)로 버스 270–504 V에서 48 V를 계약 과도 ±2 \% 안으로 조정한다. 회로에서 영전압 스위칭 실패와 역방향 경부하 이득 상승을 찾아, 자화 인덕턴스·데드타임 재설계와 경부하 두 브리지 구동으로 해결했다(TB15 5/5).
\item \textbf{R2 통합 회로(E23, TB16·TB17, 10.3절).} 실제 Stage A·조정형 Stage C를 넣은 회로에서 기본형 3/3, 창 조정 모듈 고장 강등 중 48 V 공급 2/2가 통과했다. 강등 중 버스가 380–504 V로 움직이는 동안 48 V는 ±0.03 \% 안이었다. 그동안 설계에 쓴 Stage A·Stage C 등가가 회로와 맞았다.
\item \textbf{Stage A 경부하 위상 셰딩(E24, TB18·TB19, 11장).} 통합 회로에서 찾은 경부하 영전류 스위칭의 여유 부족(1008 V 강등 뒤 10 A에서 턴오프 전류/피크 5.4 \%)을 위상 셰딩으로 해결했다(정상 운전점 26/26, 회로 5/5·2/2).
\item \textbf{충전 중 정지 협조(E26, TB20, 12장).} 충전 중 배터리 경로가 끊기면 V\textsubscript{A}가 0.01–1.1 ms 만에 400 V급 부품 상한(470 V)에 닿는다. 정지 순서와 V\textsubscript{A} 클램프(468/463 V, 1 Ω, 최대 7 kJ)를 정했고, 회로에서 4사례 중 3사례가 통과했다. 불통과 1건은 영교차 게이트가 전류 방향이 바뀌는 반주기에서 1.01 A로 꺼져 기준 1.0 A를 넘은 것이다.
\item \textbf{선행연구·산업 방식 대비 위치(E25, 13장).} 400 V급 보조부품을 쓰는 조건에서 R2는 모터·인버터 승압 충전 구조보다 보조 경로 손실이 26 \% 작고(85.0 대 114.6 W) 충전 손실도 작다. 800 V급 보조부품 설계보다는 보조 경로 손실이 15 W 크다. 따라서 R2의 이점은 400 V급 부품을 재사용하는 조건으로 한정된다. 2:1 공진 충전과 고정비 + 부분전력 조정의 원리는 선행이 있고, 같은 2:1 단의 주행 중 계속 사용과 팩 등급 조건부 창 조정·고장 강등의 조합은 조사 범위에서 찾지 못했다.
\end{itemize}

가설별 현황은 다음과 같다.

\begin{itemize}
\item H1(정격 분리)·H3(충전 중 보조버스)은 충족했다. H3의 조건은 충전전류 증가 램프 ≥ 100 µs이고, 충전 중 정지는 V\textsubscript{A} 클램프를 전제로 한다.
\item H2(강건 제어)는 평균모델·독립 구현·스위칭 모델·회로 모델에서 충족했고, 실제 Stage A·Stage C를 넣은 R2 통합 회로에서도 유지되었다.
\item H4(부하 승인)는 축약 모델에서 충족했다.
\item H5(기준 구조 대비 이점)는 사양 팩에서 기각되었고, 팩 등급별 적용 범위(7장)와 산업 방식 대비 성립 조건(13장)으로 대체했다.
\end{itemize}

남은 기간에는 검증을 먼저 끝내고 최종 보고서는 그 결과를 넣어 마지막에 쓴다(16장). 남은 검증은 영교차 게이트의 방향 반전 처리, 선행기술 전수 검색, Stage A 4모듈 경부하 운전과 부품 공차, 부하 승인 계층(SALS)의 분리·감발 연동, 부품 정격 표(V\textsubscript{A} 클램프, 제동 저항, 바이패스 소자의 선형 영역, 자성체 포화 전류, Stage C 데드타임)의 순서다.
